
While the previous section described how to construct useful run-time-monitorable properties and executable actions, guardrails must ultimately satisfy a higher-level performance objective (\perfsafety in our workflow; \Cref{fig:framework-overview}). This objective may itself be hard or inefficient to monitor at run time because it may depend on oracle knowledge, such as future accesses, for example. Moreover, multiple guardrails may be installed simultaneously, so the goal is not merely for each guardrail to behave reasonably in isolation, but for the set of installed guardrails to satisfy \perfsafety jointly. We capture this through a simple two-step verification framework. First, for each guardrail in isolation, we verify that its run-time property and action establish an appropriate local guarantee under a model of policy--system interaction. Then, under the same model, we verify that these local guarantees together imply the desired high-level objective. We next introduce the notation and semantics needed to formalize guardrails, and then present the verification framework.

\eric{
add global properties, and talk about how to specify them
}


\subsection{Semantics of Guardrails}
\label{subsec:semantics}

% \paragraph{Notation.}
We consider the application of a guardrail to a policy $\policy$ in system $S$ (where the system could be the OS as a whole, or a specific subsystem that the policy under evaluation is a part of).
Let $\states$ denote the set of possible system states for $S$ (including various statistics and state of in-kernel objects), and let $\policies$ denote the set of policies that may be deployed in the system.
% Let $\values$ denote a set of values, and let $\argspace \triangleq \values^*$ denote the set of finite lists over $\values$.
Further, each policy $\policy \in \policies$ may be instantiated with some tunable hyper-parameters $\args$, and we write the resulting deployed policy as $\ppol{\policy}{\args} \in \policies$, \aditya{something is broken. args is missing here} though notationally we often omit the subscript when it is not necessary to explicitly pull out the arguments.
For example, for hugepage promotion, Linux and CBMM are different policies, while $\ppol{\mathrm{CBMM}}{(c,b)}$ and $\ppol{\mathrm{CBMM}}{(c',b')}$ \aditya{something is broken. args are missing here} are different parameterizations of the same policy with different cost estimations $c$ and $c'$, and different benefit estimations $b$ and $b'$.

\aditya{args don't appear below! lots of brokenness}
\paragraph{Guardrail definition and semantics.}
A guardrail $G$ in a system $S$ (as defined above) is a triple $G = (\property{},\action{},\enabled)$ where $\property{}:\states\rightarrow\{0,1\}$ is the run-time monitorable condition that evaluates to 0 (false) or 1 (true) for a given state $\sysstate \in \states$, and $\action{}:\states\times\policies\rightarrow\policies$ is the specified corrective action.
The final component, $\enabled : \policies \to \{0,1\}$ identifies the policies to which the guardrail applies; thus, a guardrail only monitors a deployed policy $\ppol{\policy}{\args}$ if $\enabled(\ppol{\policy}{\args}) = 1$ (see \textsc{Guard-Disabled} rule in the semantics below).
If the evaluation of an enabled guardrail condition over the current system state $\property{}(\sysstate)$ is true, the guardrail leaves the deployed policy unchanged (see \textsc{Guard-OK} below).
Otherwise, it runs the corrective action $\action{G}$ to produce a repaired parameterized policy (see \textsc{Guard-Violation} below).
Formally, we write $\guard{G}{\sysstate}{\ppol{\policy}{\args}}{\ppol{\policy'}{\args'}}$ to denote that guardrail $G$ observes state $\sysstate \in \states$ under deployed policy $\ppol{\policy}{\args}$ and updates it to $\ppol{\policy'}{\args'}$.
The semantics are:
\begin{mathpar}
\inferrule[Guard-Disabled]{
    \enabled(\policy_\args) = 0
}{
  \guard{(\property{},\action{},\enabled)}{\sysstate}{\ppol{\policy}{\args}}{\ppol{\policy}{\args}}
}\quad
\inferrule[Guard-OK]{ 
    \enabled(\policy_\args) = 1 =
    \property{}(\sysstate)
}{
  \guard{(\property{},\action{},\enabled)}{\sysstate}{\ppol{\policy}{\args}}{\ppol{\policy}{\args}}
}\\
\inferrule[Guard-Violation]{
    \enabled(\policy_\args)
    \neq 0 = \property{}(\sysstate)
}{
    \guard{(\property{},\action{},\enabled)}{\sysstate}{\ppol{\policy}{\args}}{\action{}(\sysstate,\ppol{\policy}{\args})}
}
\end{mathpar}

\paragraph{Guarded-system semantics.}
We now define the overall system semantics induced by a single guardrail $G$.
Let a system be given by the tuple $(\states,\policies,\transitions)$, where $\states$ and $\policies$ are defined as above, and $\transitions$ is the transition relation for the system. $\TransitionsVia{\sysstate}{\ppol\policy\args}{\sysstate'}$ indicates that the system steps from $\sysstate$ to $\sysstate'$ under policy $\ppol\policy\args$.
The guarded system is then, an extension of $S$, where at each step, it first checks whether the guardrail is violated and repairs the deployed parameterized policy if needed.
This is captured by the rule:
\begin{mathpar}
\inferrule[Guarded-Transition]{
    \TransitionsVia{\sysstate}{\ppol{\policy}{\args'}}{\sysstate'}
    \quad
    \guard{G}{{\sysstate'}}{\ppol{\policy}{\args}}{\ppol{\policy'}{\args'}} 
}{
    \transition{G}{\runtimestate{\sysstate}{\ppol{\policy}{\args}}}{\runtimestate{\sysstate'}{\ppol{\policy'}{\args'}}}
}
\end{mathpar}
Informally, the above rule implies the following guarded system semantics: detect whether the current state violates the guardrail for the currently deployed policy, repair the deployed parameterized policy if needed, and then take the next system step under the resulting policy. \aditya{this sentence does not parse and seems like an overclaim} The semantics of multiple guardrails can be defined similarly by guardrails in some nondeterministic order.


\subsection{Formal Verification of Performance Properties}
\label{subsec:verification-conditions}

\ds{Motivating the \perfsafety property with examples.}
The \textsc{Guarded-Transition} rule above induces a transition system over runtime states.
The verification question is whether the local guarantees attached to the guardrails are strong enough to characterize safe states of this transition system, and whether those local guarantees for multiple guardrails together imply the desired \perfsafety property.
% Our verification framework below reasons over an abstract model $M$ that over-approximates $\transitions$: any transition possible under $\transitions$ is also possible under $M$, but $M$ may permit additional transitions that simplify the proof.

The guardrail components $\property{G}$, $\action{G}$, and $\policies_G$ are already defined in \Cref{subsec:semantics}.
The verification framework introduces additional predicates per guardrail: a \emph{local guarantee} $\varphi(\sysstate,\ppol{\policy}{\args})$, which is a predicate that is guaranteed to hold both after the guardrail is triggered, and when the guardrail is not triggered. Additionally, we have a \emph{global \perfsafety assertion} $\Phi(\sysstate,\ppol{\policy}{\args})$, which is the high-level objective to be established.
It uses an abstract model $M(\sysstate,\ppol{\policy}{\args},\sysstate')$ in place of the concrete transition relation $\transitions$; $M$ over-approximates $\transitions$ so any transition possible under $\transitions$ is also possible under $M$. Consequently, if a set of guardrails satisfies $\Phi(\sysstate, \ppol{\policy}{\args})$ under model $M$, it is also safe w.r.t. the transition system $\transitions$.

\paragraph{Local Verification.}
For each guardrail $G = (\property{}, \action{}, \enabled)$ and property $\varphi$, the verifier checks the following:
\[
\begin{array}{c@{\;\;}l}
M(\sysstate, \policy, \sysstate') \wedge \enabled(\policy) \wedge \property{}(\sysstate') \Rightarrow \varphi(\policy, \sysstate') & \textsc{Maintain} \\
M(\sysstate,\policy, \sysstate') \wedge \enabled(\policy) \wedge \neg P(\sysstate') 
\Rightarrow \varphi(\action{}(\sysstate, \policy), \sysstate')
& 
\textsc{Recover}
\end{array}
\]


% \[
% \begin{array}{r@{\;\;\;\;}l}
% \textsc{Ok} & 
% \bigl( (\ppol{\policy}{\args} \in \policies_G) \land \property{G}(\sysstate) \Rightarrow \varphi(\sysstate',\ppol{\policy}{\args})\bigr)\\
% \textsc{Viol} & \bigl((\ppol{\policy}{\args} \in \policies_G) \land \neg\property{G}(\sysstate) \Rightarrow \varphi(\sysstate',\action{G}(\sysstate,\ppol{\policy}{\args}))\bigr) \\[0.4em]
% \textsc{Global} & I(\sysstate,\ppol{\policy}{\args}) \Rightarrow \Phi(\sysstate,\ppol{\policy}{\args})
% \end{array}
% \]
The \textsc{Maintain} condition says that, under system model $M$, if a guardrail is enabled for a policy $\enabled(\policy)$, and the current state already satisfies the guardrail condition $\property{}(\sysstate')$, the local guarantee $\varphi$ holds.
The \textsc{Recover} condition covers the complementary case: when the guardrail applies ($\enabled(\policy)$) but $\property{}$ does not hold, the corrective action $\action{}$ establishes $\varphi$.

\paragraph{Performance Safety Verification}
Now, to reason globally, we need to ensure that the system $M$, monitored by a set of guardrails $\{G_1, \ldots, G_n\}$ \aditya{proves or ensures?} proves ensures $\Phi(\sysstate, \policy)$ always holds. To do this we construct a formula $I$ that is the invariant of the combined system:
\[
I =
M(\sysstate, \policy, \sysstate') \wedge
\bigwedge_{i=1}^n \textsc{Maintain}_i(\sysstate, \policy, \sysstate') \wedge 
\textsc{Recover}_i(\sysstate,\policy, \sysstate')
\]
where $\textsc{Maintain}_i$ and $\textsc{Recover}_i$ for $i = 1,\ldots, n$ are the conditions above instantiated for each of the $n$ guardrails. Then, aside from showing that the the initial configuration of the system satisfies the constraints, we need to prove only that this invariant implies the performance safety condition $\Phi$
\[
\begin{array}{c@{\qquad}l}
    I(\sysstate, \policy, \sysstate) \Rightarrow \Phi(\policy, \sysstate)
    & \text{\perfsafety}
\end{array}  
\]
which says that the set of local guarantees applies  the system-wide assertion $\Phi$. Together, these three checks say that the local conditions are collectively sufficient.

\paragraph{Soundness.} In order for this approach to be sound, we need to deal with some concurrency effects. Because guardrails trigger nondeterministically, we need to ensure that they do not interfere with each others' guarantees. This is captured in the following soundness condition:
\[
\inferrule[Sound(\ensuremath{\pi},\ensuremath{i},\ensuremath{j})]{}
{\qquad\enabled_i(\policy) \wedge 
\enabled_j(\policy) \wedge
\varphi_j(\policy, \sysstate) \wedge
\neg \property{i}(\policy, \sysstate) \Rightarrow
\varphi_j(\action i(\policy), \sysstate) }
\]

% \[
% \begin{array}{r@{~\Rightarrow~}l@{\qquad}l}
%    \neg \property i(\sysstate) \wedge \neg \property j (\sysstate) &
%    \varphi_i(\action j, \sysstate)
%    \wedge
%    \varphi_j(\action i, \sysstate) 
%     & \textsc{Both}_{i,j}
%     \\
%     \property i(\sysstate) \wedge \neg \property j(\sysstate)
%     &
%     \varphi_i(\action j, \sysstate).
%     & \textsc{One}_{i,j}
% \end{array}  
% \]
which must hold (in the system model) for each policy $\policy$ and each pair \aditya{of what?} in the set of guardrails $\{G_1, \ldots, G_n\}$. It says that if two guardrails are enabled on the same policy, they must each ensure each other's guarantees. The verification procedure is sound under this condition.

% If the initial runtime state satisfies $I$, if \textsc{Ok} and \textsc{Viol} hold for each guardrail under its model $M$ and chosen $\varphi$, and if \textsc{Global} holds, then every runtime state reachable under the guarded transition relation satisfies $\Phi$. In other words, the conjunction of the local guardrail guarantees is an inductive invariant for the transition system generated from the system model $M$ and that invariant entails the global \perfsafety assertion.


\iffalse

% <=================== FULL FORMAL NOTATION ==================>

\subsection{Semantics of Guardrails}
\label{subsec:semantics}

\paragraph{Notation.}
% A guardrail can apply to the execution of the OS under some policy of interest.
We consider the application of a guardrail to a policy $\policy$ in system $S$ (where the system could be the OS as a whole, or a specific subsystem that the policy under evaluation is a part of).
Let $\values$ denote a set of values, and let $\argspace \triangleq \values^*$ denote the set of finite lists over $\values$.
Each policy $\policy \in \policies$ may be instantiated with arguments $\args \in \argspace$, and we write the resulting deployed policy as $\ppol{\policy}{\args}$.
We write
\(
\ParametricPolicies{\policies}{\argspace} \triangleq \{\, \ppol{\policy}{\args} \mid \policy \in \policies,\ \args \in \argspace \,\}
\)
for the set of parameterized policies induced by $\policies$ and argument lists from $\argspace$.
Here, `policy' is used in the usual systems sense: Linux and CBMM are different policies, while $\ppol{\mathrm{CBMM}}{(c,b)}$ and $\ppol{\mathrm{CBMM}}{(c',b')}$ are different parameterizations of the same policy with different cost estimations $c$ and $c'$, and different benefit estimations $b$ and $b'$.
Let $\states$ denote the set of possible system states for $S$ (including various statistics and state of in-kernel objects), and let $\policies$ denote the set of policies that may be deployed in the system.
% We also write
% \[
% \ActionSet \triangleq \states \times \ParametricPolicies{\policies}{\argspace} \xrightarrow{} \ParametricPolicies{\policies}{\argspace}
% \]
% At runtime, the system runs a parameterized policy $\ppol{\policy}{\args} \in \ParametricPolicies{\policies}{\argspace}$.
% We will now formally describe how a guardrail operates in this system.

\paragraph{Guardrail definition and semantics.}
A guardrail $G \in \guardrails$ in a system $S$ (as defined above) is given by the triple $(\policies_G,\states_G,\action{G})$ where $\policies_G \subseteq \policies$ is the set of policies monitored by the guardrail, $\states_G \subseteq \states$ is the set of states that the guardrail considers acceptable, and $\action{G} \in \ActionSet$.
Informally, the guardrail monitors the current parameterized policy $\ppol{\policy}{\args}$ only when its policy $\policy$ belongs to $\policies_G$.
If the current system state $\sysstate$ lies in $\states_G$, the guardrail leaves the deployed policy unchanged.
Otherwise, it runs the corrective action $\action{G}$ to produce a repaired parameterized policy.
Formally, we write $\guard{G}{\sysstate}{\ppol{\policy}{\args}}{\ppol{\policy'}{\args'}}$ to denote that guardrail $G$ observes state $\sysstate \in \states$ under deployed policy $\ppol{\policy}{\args}$ and updates it to $\ppol{\policy'}{\args'}$.
The semantics are:
\begin{mathpar}
\inferrule[Guard-Unenabled]{
    \policy \notin \policies_G
}{
  \guard{(\policies_G,\states_G,\action{G})}{\sysstate}{\ppol{\policy}{\args}}{\ppol{\policy}{\args}}
}\quad
\inferrule[Guard-OK]{ 
    \policy \in \policies_G \quad
    \sysstate \in \states_G
}{
  \guard{(\policies_G,\states_G,\action{G})}{\sysstate}{\ppol{\policy}{\args}}{\ppol{\policy}{\args}}
}\\
\inferrule[Guard-Violation]{
    \policy \in \policies_G \quad
    \sysstate \notin \states_G
}{
    \guard{(\policies_G,\states_G,\action{G})}{\sysstate}{\ppol{\policy}{\args}}{\action{G}(\sysstate,\ppol{\policy}{\args})}
}
\end{mathpar}

\paragraph{Guarded-system semantics.}
We now define the outer system semantics induced by a single guardrail $G$.
Let a system be given by the tuple $(\states,\policies,\transitions)$, where $\sysstate, \policy \transitions \sysstate'$ for $\sysstate',\sysstate \in \states$ and $\policy \in \policies$ indicates that the system may step from $\sysstate$ to $\sysstate'$ under policy $\policy$.
The guarded system is an extension of the underlying system, where at each step, it first checks whether the guardrail is violated and repairs the deployed parameterized policy if needed.
This is captured by the rule:
\begin{mathpar}
\inferrule[Guarded-Transition]{
    \guard{G}{\sysstate}{\ppol{\policy}{\args}}{\ppol{\policy'}{\args'}} \quad
    \sysstate,\ppol{\policy'}{\args'} \transitions \sysstate'
}{
    \transition{G}{\runtimestate{\sysstate}{\ppol{\policy}{\args}}}{\runtimestate{\sysstate'}{\ppol{\policy'}{\args'}}}
}
\end{mathpar}
Thus, the guarded system semantics is simple: detect whether the current state violates the guardrail for the currently deployed policy, repair the deployed parameterized policy if needed, and then take the next system step under the resulting policy. The semantics of multiple guardrails can be defined similarly by applying the above rule for each guardrail in some fixed order.


\subsection{Formal Verification of Performance Properties}
\label{subsec:verification-conditions}

% Performance Properties are safety properties
% we can verify them just like any other safety properties.
% Though we require an abstraction of the system, and a way to represent actions.

The guarded transition rules above induce an outer transition system over runtime states.
The verification question is whether the local guarantees attached to the guardrails are strong enough to characterize safe states of this transition system, and whether those local guarantees together imply the desired system-wide property.
For this purpose, we summarize a fixed guardrail using the following symbolic interface:
\[
\begin{array}{r@{\;\;:\;\;}l}
M(\sysstate,\ppol{\policy}{\args},\sysstate') & \text{system model step relation} \\[0.2em]
P(\ppol{\policy}{\args}) & \text{policy applicability predicate} \\[0.2em]
S(\sysstate) & \text{guardrail condition} \\[0.2em]
A(\sysstate,\ppol{\policy}{\args}) & \text{corrective policy update} \\[0.2em]
\varphi(\sysstate,\ppol{\policy}{\args}) & \text{local guarantee} \\[0.2em]
\Phi(\sysstate,\ppol{\policy}{\args}) & \text{global \perfsafety assertion}
\end{array}
\]
Here $M$ is simply the system transition relation from \Cref{subsec:formal-guardrails}, written as a predicate.
Likewise, $P$, $S$, $A$, and $\varphi$ are the symbolic views of one guardrail's applicability condition, state condition, corrective action, and local guarantee.

\paragraph{Verification interface.}
Under the corresponding system-model step $M$, the verifier checks the following three conditions; free variables are implicitly universally quantified:
\[
\begin{array}{r@{\;\;\;\;}l}
\textsc{Ok} & 
\bigl(P(\ppol{\policy}{\args}) \land S(\sysstate) \Rightarrow \varphi(\sysstate',\ppol{\policy}{\args})\bigr)\\
\textsc{Viol} & \bigl(P(\ppol{\policy}{\args}) \land \neg S(\sysstate) \Rightarrow \varphi(\sysstate',A(\sysstate,\ppol{\policy}{\args}))\bigr) \\[0.4em]
\textsc{Global} & I(\sysstate,\ppol{\policy}{\args}) \Rightarrow \Phi(\sysstate,\ppol{\policy}{\args})
\end{array}
\]
where $I$ is the conjunction of \textsc{OK} and $\textsc{Viol}$ for all guardrails.

The \textsc{Ok} condition says that, when the deployed policy is one to which the guardrail applies and the current state already satisfies the state condition $S$, one system-model step must land in a runtime state satisfying the local guarantee $\varphi$. The \textsc{Viol} condition covers the complementary case: when the guardrail applies but the current state does not satisfy $S$, applying the corrective update $A$ and then taking one system-model step must still establish $\varphi$. Finally, \textsc{Global} says that if the conjunction $I$ of all local guarantees holds at a runtime state, then the system-wide assertion $\Phi$ must also hold there. Together, these three checks say that the local conditions are strong enough individually and sufficient collectively.

\paragraph{Soundness.}
If the initial runtime state satisfies $I$, if \textsc{Ok} and \textsc{Viol} hold for each guardrail under its chosen predicates $P$, $S$, $A$, and $\varphi$, and if \textsc{Global} holds, then every runtime state reachable under the guarded transition relation satisfies $\Phi$.
In other words, the conjunction of the local guardrail guarantees is an inductive invariant for the outer transition system generated from the system model, and that invariant entails the global \perfsafety assertion.

\fi